\documentclass[11pt,letterpaper]{article}

\usepackage{geometry}
\geometry{margin=1in}

\usepackage{graphicx}
\usepackage{booktabs}
\usepackage{amsmath}
\usepackage{amssymb}
\usepackage{hyperref}
\usepackage{xcolor}
\usepackage{float}
\usepackage{caption}
\usepackage{subcaption}
\usepackage{multirow}
\usepackage{array}
\usepackage{colortbl}
\usepackage{algorithm}
\usepackage{algorithmic}

\definecolor{bestcolor}{RGB}{200,255,200}
\definecolor{goodcolor}{RGB}{255,255,200}

\title{\textbf{Thermal Control Algorithms for\\Primary Mirror Temperature Matching}\\[0.5em]
\large Report 2: Rate-Matching Control for Rubin Observatory}

\author{Christopher Stubbs and Johnny Esteves\\[0.5em]
\small\textit{Analysis performed with extensive use of Claude Code}}
\date{January 2026}

\begin{document}

\maketitle

\begin{abstract}
This report presents control algorithms for matching the Rubin Observatory primary mirror temperature to ambient conditions throughout the night. We analyze a \textbf{dual-setpoint control} approach where the dome interior HVAC maintains a \emph{constant} temperature before sunset, and the mirror naturally equilibrates to this interior temperature. At sunset, the dome opens and interior air instantly equilibrates to exterior ambient, while the mirror (with $\tau = 2$ hour time constant) is then actively controlled with rate-matching.

Simulations across 769 nights demonstrate that the optimal strategy uses:
\begin{itemize}
    \item \textbf{Pre-sunset}: Dome HVAC holds interior at predicted sunset temperature minus 0.3°C (FLAT, not tracking exterior)
    \item \textbf{Overnight}: Rate-matching control with 0.3°C cold bias
\end{itemize}

This achieves \textbf{0.44°C RMS} tracking error overnight, with \textbf{68\%} of observing time in the ideal range (mirror 0 to 0.5°C cooler than ambient) and only \textbf{15\%} of time with the mirror warmer than ambient.
\end{abstract}

\tableofcontents

\newpage

%=============================================================================
\section{Introduction}
%=============================================================================

\subsection{The Control Problem}

The primary mirror thermal control problem has two distinct phases:

\begin{enumerate}
    \item \textbf{Pre-sunset Phase}: Drive mirror to predicted sunset conditions
    \item \textbf{Overnight Phase}: Track ambient temperature as it evolves
\end{enumerate}

Traditional approaches focus only on matching temperature. This report demonstrates that matching \emph{both} temperature and rate of change yields significantly better overnight performance.

\subsection{Physical Model}

The system has two controllable temperatures:

\begin{enumerate}
    \item \textbf{Dome interior air temperature} $T_\text{dome}$: Controlled by HVAC with fast response ($\tau_\text{dome} \approx 30$ minutes)
    \item \textbf{Mirror temperature setpoint} $T_\text{mirror,sp}$: Active heating/cooling system
\end{enumerate}

The mirror temperature follows first-order dynamics:
\begin{equation}
    \frac{dT_\text{mirror}}{dt} = \frac{T_\text{setpoint} - T_\text{mirror}}{\tau}
    \label{eq:dynamics}
\end{equation}

where $\tau = 2$ hours is the thermal time constant. The setpoint rate of change is limited to $\pm 2°C/\text{hour}$ to prevent thermal gradients.

\textbf{Key operational constraint}: Before sunset, the mirror exchanges heat with the dome interior air (not exterior ambient). At sunset, the dome opens and interior air \emph{immediately} equilibrates to exterior ambient temperature.

\subsection{Why Rate Matters}

If the mirror has the correct temperature but wrong rate at sunset:
\begin{itemize}
    \item The mirror will drift away from ambient as the night progresses
    \item Temperature errors accumulate, degrading image quality
    \item Recovery requires aggressive control that may cause thermal gradients
\end{itemize}

Matching both $T$ and $\dot{T}$ at sunset ensures the mirror ``tracks'' ambient naturally through the night.

%=============================================================================
\section{Control Strategy Overview}
%=============================================================================

\subsection{Phase 1: Pre-Sunset Control}

Before sunset, the dome is closed and the interior environment is controlled by HVAC. The control strategy is simple:

\begin{equation}
    \boxed{T_\text{HVAC} = T_\text{pred} - 0.3°C \quad \text{(constant)}}
    \label{eq:presunset}
\end{equation}

where $T_\text{pred}$ is the predicted sunset temperature from the Random Forest model.

\textbf{Key points}:
\begin{itemize}
    \item The HVAC setpoint is \textbf{FLAT} (constant)---it does \emph{not} track exterior temperature
    \item The mirror naturally equilibrates to the dome interior temperature
    \item No active mirror temperature control is needed before sunset
    \item The 0.3°C cold bias ensures the mirror starts slightly cold at sunset
\end{itemize}

The dome air equilibrates quickly ($\tau_\text{dome} = 30$ min), and the mirror equilibrates to the dome interior over its 2-hour time constant. Starting 4+ hours before sunset ensures both are at the target temperature.

\subsection{Phase 2: Sunset Transition}

At sunset, the dome opens and:
\begin{itemize}
    \item Dome interior air \textbf{immediately} equilibrates to exterior ambient
    \item Mirror is at approximately $T_\text{sunset} - 0.4°C$ (slightly cold, as desired)
    \item Control switches to overnight tracking mode
\end{itemize}

\subsection{Phase 3: Overnight Tracking}

After sunset, the algorithm uses rate-matching to track the evolving ambient:
\begin{equation}
    \boxed{T_\text{setpoint} = T_\text{ambient} + \tau \cdot \dot{T}_\text{ambient} - \sigma_\text{night}}
    \label{eq:overnight}
\end{equation}

where:
\begin{itemize}
    \item $T_\text{ambient}$: current exterior ambient temperature (measured)
    \item $\dot{T}_\text{ambient}$: ambient cooling rate computed from 1-hour history
    \item $\tau = 2$ hours: mirror thermal time constant
    \item $\sigma_\text{night} = 0.3°C$: overnight cold bias (same as pre-sunset)
\end{itemize}

The overnight rate is estimated as:
\begin{equation}
    \dot{T}_\text{ambient} = \frac{T_\text{now} - T_{-1\text{h}}}{1\text{ hour}}
\end{equation}

\textbf{Physical Interpretation}: If ambient is cooling at 0.5°C/hour, the setpoint is $2 \times 0.5 = 1°C$ colder than ambient. This causes the mirror to warm toward the setpoint at the same rate that ambient is cooling, achieving rate-matching.

%=============================================================================
\section{Algorithm Pseudocode}
%=============================================================================

The complete control algorithm is presented below:

\begin{algorithm}[H]
\caption{Rate-Matching Thermal Control}
\label{alg:control}
\begin{algorithmic}[1]
\STATE \textbf{Constants:}
\STATE $\tau \gets 2.0$ hours \COMMENT{thermal time constant}
\STATE $\sigma_\text{pre} \gets 0.20$°C \COMMENT{pre-sunset cold bias}
\STATE $\sigma_\text{night} \gets 0.15$°C \COMMENT{overnight cold bias}
\STATE $\Delta t \gets 15$ minutes \COMMENT{control update interval}
\STATE $R_\text{max} \gets 2.0$°C/hour \COMMENT{max setpoint rate of change}

\STATE
\STATE \textbf{Pre-sunset Phase} (4h to sunset):
\WHILE{time $<$ sunset}
    \STATE $T_\text{pred}, \dot{T}_\text{pred} \gets$ \textsc{PredictSunset}(current\_time)
    \STATE $\dot{T}_\text{ambient} \gets$ \textsc{MeasureAmbientRate}()
    \STATE $t_\text{remain} \gets$ hours until sunset
    \STATE $\beta \gets \max(0, 1 - t_\text{remain}/1.5)$ \COMMENT{blend factor}
    \STATE $\dot{T}_\text{eff} \gets (1-\beta) \cdot \dot{T}_\text{pred} + \beta \cdot \dot{T}_\text{ambient}$
    \STATE $T_\text{raw} \gets T_\text{pred} + \tau \cdot \dot{T}_\text{eff} - \sigma_\text{pre}$
    \STATE $T_\text{setpoint} \gets$ \textsc{RateLimit}($T_\text{raw}$, $T_\text{prev}$, $R_\text{max}$, $\Delta t$)
    \STATE \textsc{ApplySetpoint}($T_\text{setpoint}$)
    \STATE Wait $\Delta t$
\ENDWHILE

\STATE
\STATE \textbf{Overnight Phase} (sunset to sunrise):
\STATE $T_\text{prev} \gets T_\text{ambient}$ \COMMENT{initialize rate calculation}
\WHILE{time $<$ sunrise}
    \STATE $T_\text{now} \gets$ \textsc{MeasureAmbient}()
    \STATE $\dot{T}_\text{ambient} \gets (T_\text{now} - T_\text{prev}) / \Delta t_\text{rate}$
    \STATE $T_\text{raw} \gets T_\text{now} + \tau \cdot \dot{T}_\text{ambient} - \sigma_\text{night}$
    \STATE $T_\text{setpoint} \gets$ \textsc{RateLimit}($T_\text{raw}$, $T_\text{prev}$, $R_\text{max}$, $\Delta t$)
    \STATE \textsc{ApplySetpoint}($T_\text{setpoint}$)
    \STATE $T_\text{prev} \gets T_\text{now}$
    \STATE Wait $\Delta t$
\ENDWHILE

\STATE
\STATE \textbf{RateLimit}($T_\text{raw}$, $T_\text{prev}$, $R_\text{max}$, $\Delta t$):
\STATE \quad $\Delta T_\text{max} \gets R_\text{max} \cdot \Delta t$
\STATE \quad \textbf{return} $\text{clip}(T_\text{raw}, T_\text{prev} - \Delta T_\text{max}, T_\text{prev} + \Delta T_\text{max})$
\end{algorithmic}
\end{algorithm}

\subsection{Implementation Notes}

\begin{enumerate}
    \item \textbf{Setpoint rate limiting}: The setpoint rate of change is limited to $\pm 2°C$/hour to prevent thermal gradients across the mirror. This is critical---rapid setpoint changes would create temperature differentials that degrade image quality.
    \begin{equation}
        |T_\text{setpoint}(t) - T_\text{setpoint}(t-\Delta t)| \leq R_\text{max} \cdot \Delta t = 0.5°C \text{ per 15 min}
    \end{equation}

    \item \textbf{Rate smoothing}: Use exponential moving average for overnight rate to reduce noise:
    \begin{equation}
        \dot{T}_\text{smooth} = \alpha \cdot \dot{T}_\text{new} + (1-\alpha) \cdot \dot{T}_\text{prev}, \quad \alpha = 0.3
    \end{equation}

    \item \textbf{Ambient rate estimation}: Compute rate from 1-hour temperature history to smooth out short-term fluctuations while remaining responsive to actual trends.

    \item \textbf{Prediction updates}: During pre-sunset phase, predictions improve as sunset approaches. Update every 15 minutes using the latest forecast.

    \item \textbf{Pre-sunset to overnight transition}: The dome environment differs from outdoor ambient during the day, and HVAC energy limits prevent tracking outdoor ambient before sunset. Therefore:
    \begin{itemize}
        \item \textbf{Early pre-sunset} ($>$1.5h before sunset): Use predicted sunset rate ($\dot{T}_\text{pred}$) in the setpoint formula
        \item \textbf{Near sunset} (1.5h to 0h): Linearly blend from predicted rate to measured ambient rate
        \item \textbf{At/after sunset}: Use measured ambient rate exclusively
    \end{itemize}
    The blending function is:
    \begin{equation}
        \dot{T}_\text{effective} = (1-\beta) \cdot \dot{T}_\text{pred} + \beta \cdot \dot{T}_\text{ambient}
    \end{equation}
    where $\beta = \max(0, 1 - t_\text{to\_sunset}/1.5\text{h})$. This ensures smooth setpoint transition as the dome opens at sunset.
\end{enumerate}

%=============================================================================
\section{Simulation Results}
%=============================================================================

\subsection{Test Methodology}

We simulated both temperature-only and rate-matching control across 770 nights of historical data:
\begin{itemize}
    \item Pre-sunset phase: 4 hours with RF-based predictions
    \item Overnight phase: Sunset to sunrise with measured ambient tracking
    \item Mirror dynamics: First-order response with $\tau = 2$ hours
\end{itemize}

\subsection{At-Sunset Performance}

\begin{table}[H]
\centering
\caption{Performance Comparison at Sunset}
\label{tab:sunset}
\begin{tabular}{l|cc}
\toprule
\textbf{Metric} & \textbf{Temperature-Only} & \textbf{Rate-Matching} \\
\midrule
Temperature RMS & 0.80°C & 0.69°C \\
Temperature Bias & +0.27°C (warm) & $-0.18$°C (cold) \\
Rate RMS & 0.87°C/h & 0.34°C/h \\
Rate Bias & +0.65°C/h & +0.04°C/h \\
\% Days Warm & 64.7\% & 39.9\% \\
\bottomrule
\end{tabular}
\end{table}

\textbf{Key Results}:
\begin{itemize}
    \item Rate-matching reduces rate error by 61\%
    \item Temperature error also improves by 13\%
    \item Cold bias shifts from 65\% warm to 60\% cold (as desired)
\end{itemize}

\subsection{Overnight Performance}

\begin{table}[H]
\centering
\caption{Overnight Tracking Performance (Sunset to Sunrise)}
\label{tab:overnight}
\begin{tabular}{l|cc}
\toprule
\textbf{Metric} & \textbf{Temperature-Only} & \textbf{Rate-Matching} \\
\midrule
\rowcolor{bestcolor}
Temperature RMS & 0.53°C & \textbf{0.30°C} \\
Rate RMS & 0.26°C/h & 0.14°C/h \\
Improvement & --- & 44\% \\
\bottomrule
\end{tabular}
\end{table}

Rate-matching control achieves \textbf{44\% reduction} in overnight tracking error. This is the primary benefit: better initial conditions propagate through the entire night.

\subsection{Cold Bias Optimization}

We systematically tested cold bias values from 0.0°C to 0.4°C to find the optimal operating point. The goal is to keep the mirror slightly cooler than ambient (0 to 0.5°C) while minimizing RMS tracking error.

\begin{table}[H]
\centering
\caption{Cold Bias Optimization Results (769 nights, both biases equal)}
\label{tab:cold_bias}
\begin{tabular}{c|cccc}
\toprule
\textbf{Cold Bias} & \textbf{RMS Error} & \textbf{Mean Offset} & \textbf{\% Warm} & \textbf{\% Ideal} \\
\midrule
0.0°C & 0.33°C & +0.04°C & 55\% & 40\% \\
0.1°C & 0.34°C & $-0.06$°C & 35\% & 58\% \\
0.2°C & 0.38°C & $-0.16$°C & 22\% & 67\% \\
\rowcolor{bestcolor}
0.3°C & 0.44°C & $-0.26$°C & 15\% & 68\% \\
0.4°C & 0.51°C & $-0.36$°C & 11\% & 62\% \\
\bottomrule
\end{tabular}
\end{table}

The \textbf{0.3°C cold bias} (for both pre-sunset and overnight) provides the best operational balance:
\begin{itemize}
    \item Reasonable RMS error (0.44°C)
    \item Appropriate cold offset ($-0.26$°C mean, ``a few tenths cold'')
    \item Maximum time in ideal range (68\% in 0 to $-0.5$°C window)
    \item Only 15\% of time warm (greatly reduced from no-bias case)
\end{itemize}

Figure~\ref{fig:cold_bias} shows the trade-off analysis.

\begin{figure}[H]
\centering
\includegraphics[width=\textwidth]{fig76_fixed_parameter_sweep.png}
\caption{Parameter sweep for pre-sunset and overnight cold bias values. (A) RMS error is minimized at low biases but (B) \% time warm is unacceptably high. (C) \% time in ideal range is maximized around 0.3°C for both biases. (D) Mean offset shifts cold with increasing bias. The optimal operating point of 0.3°C/0.3°C balances RMS (0.44°C), cold preference (only 15\% warm), and time in ideal range (68\%).}
\label{fig:cold_bias}
\end{figure}

\subsection{Overnight Tracking Examples}

Figure~\ref{fig:overnight_examples} shows the control algorithm tracking ambient temperature through 9 representative nights.

\begin{figure}[H]
\centering
\includegraphics[width=\textwidth]{fig75_fixed_dual_setpoint_examples.png}
\caption{Dual-setpoint control examples showing the complete operational cycle. Before sunset (negative hours): the mirror (green) is FLAT at the HVAC setpoint (red line), NOT tracking exterior ambient (blue). At sunset (t=0): the dome opens and interior air instantly equilibrates to exterior ambient. Overnight: mirror tracks ambient with rate-matching control. Blue shading indicates mirror cooler than ambient (good); pink shading indicates warmer (less desirable).}
\label{fig:overnight_examples}
\end{figure}

The mirror tracks ambient closely across diverse conditions, with most of the shaded area showing blue (cold) rather than pink (warm).

\subsection{Error Distribution}

Figure~\ref{fig:error_dist} shows the distribution of mirror-minus-ambient temperature across all 770 nights.

\begin{figure}[H]
\centering
\includegraphics[width=0.8\textwidth]{fig77_fixed_error_distribution.png}
\caption{Distribution of mirror-ambient temperature offset with 0.3°C pre-sunset and overnight cold bias. The distribution is centered at $-0.26$°C (a few tenths cold, as desired), with most of the mass in the ideal range (0 to $-0.5$°C, green shading). Only 15\% of observations fall in the warm region (pink shading).}
\label{fig:error_dist}
\end{figure}

\subsection{Visualization}

Figure~\ref{fig:comparison} shows the comparison of control strategies.

\begin{figure}[H]
\centering
\includegraphics[width=\textwidth]{fig42_rate_matching_real_predictions.png}
\caption{Rate-matching vs temperature-only control using real RF predictions. Temperature-only control is shown in green, rate-matching in red. Rate-matching shows tighter error distributions and reduced warm bias across all panels.}
\label{fig:comparison}
\end{figure}

%=============================================================================
\section{Example Days}
%=============================================================================

Figure~\ref{fig:examples} shows the control algorithm operation on representative days.

\begin{figure}[H]
\centering
\includegraphics[width=\textwidth]{fig34_example_days_overview.png}
\caption{Example temperature trajectories across seasons. Blue solid: actual ambient. Magenta dashed: predicted sunset temperature (with pink uncertainty band). Orange dotted: mirror setpoint. Green solid: simulated mirror temperature. Vertical orange line marks sunset.}
\label{fig:examples}
\end{figure}

The algorithm successfully:
\begin{itemize}
    \item Tracks the evolving sunset prediction (orange line converges to actual)
    \item Applies lead compensation (mirror cooler than target when cooling expected)
    \item Achieves smooth transitions at sunset
    \item Maintains tight tracking overnight
\end{itemize}

%=============================================================================
\section{Cold Bias Strategy}
%=============================================================================

\subsection{Rationale}

A slight cold bias is preferred because:
\begin{enumerate}
    \item \textbf{Seeing quality}: Cold mirror has rising warm air (good). Warm mirror has sinking cold air (turbulence).
    \item \textbf{Condensation risk}: Warm mirror cannot condense moisture; cold mirror can be actively dried.
    \item \textbf{Asymmetric correction}: Easier to warm a cold mirror (use heaters) than cool a warm one (limited by ambient).
\end{enumerate}

\subsection{Fixed vs Adaptive Bias}

We investigated whether adapting the cold bias based on conditions could improve performance:

\begin{itemize}
    \item \textbf{Rate-adaptive}: Increase bias when cooling rapidly (higher risk of mirror lagging)
    \begin{equation}
        \sigma = \sigma_\text{base} + k \cdot |\dot{T}| \quad \text{when } \dot{T} < 0
    \end{equation}
    \item \textbf{Uncertainty-adaptive}: Increase bias when prediction uncertainty is high
    \item \textbf{Combined}: Both rate and uncertainty dependence
\end{itemize}

\begin{table}[H]
\centering
\caption{Fixed vs Adaptive Bias Comparison (769 nights)}
\label{tab:adaptive}
\begin{tabular}{l|cccc}
\toprule
\textbf{Strategy} & \textbf{RMS} & \textbf{Mean} & \textbf{\% Warm} & \textbf{\% Ideal} \\
\midrule
Fixed 0.2°C & 0.37°C & $-0.21$°C & 17\% & \textbf{72\%} \\
\rowcolor{bestcolor}
Fixed 0.3°C & 0.44°C & $-0.31$°C & \textbf{11\%} & 70\% \\
Rate-Adaptive & 0.45°C & $-0.29$°C & 15\% & 63\% \\
Combined & 0.43°C & $-0.26$°C & 15\% & 67\% \\
\bottomrule
\end{tabular}
\end{table}

\textbf{Key findings}:
\begin{itemize}
    \item Fixed 0.2°C maximizes time in ideal range (72\%)
    \item Fixed 0.3°C minimizes warm time (11\%) with slightly less ideal time (70\%)
    \item Adaptive strategies do not clearly outperform fixed bias
    \item Added complexity of adaptive schemes is not justified by performance gains
\end{itemize}

\subsection{Recommended Bias Values}

Based on operational priorities:
\begin{itemize}
    \item \textbf{If maximizing ideal range}: Use 0.2°C bias (72\% ideal, 17\% warm)
    \item \textbf{If minimizing warm time}: Use 0.3°C bias (70\% ideal, 11\% warm)
\end{itemize}

We recommend \textbf{0.3°C} as it provides strong protection against warm conditions while maintaining excellent ideal-range performance. The fixed approach is preferred for simplicity and robustness.

\subsection{Results Summary}

\begin{table}[H]
\centering
\caption{Temperature Bias Analysis}
\label{tab:bias}
\begin{tabular}{l|cc}
\toprule
\textbf{Control Strategy} & \textbf{\% Days Warm} & \textbf{Mean Bias} \\
\midrule
Temperature-Only & 64.7\% & +0.27°C \\
\rowcolor{bestcolor}
Rate-Matching (0.3°C bias) & 11\% & $-0.31$°C \\
\midrule
\textit{Target} & \textit{$<$20\%} & \textit{$<$0°C} \\
\bottomrule
\end{tabular}
\end{table}

Rate-matching achieves the desired cold bias: 60\% of nights are slightly cold rather than warm.

%=============================================================================
\section{Robustness Analysis}
%=============================================================================

\subsection{Sensitivity to Time Constant}

The algorithm uses $\tau = 2$ hours. What if the actual time constant differs?

\begin{table}[H]
\centering
\caption{Sensitivity to Time Constant Mismatch}
\label{tab:tau}
\begin{tabular}{l|ccc}
\toprule
\textbf{Actual $\tau$} & \textbf{1.5 h} & \textbf{2.0 h} & \textbf{2.5 h} \\
\midrule
Algorithm uses $\tau=2$h & & & \\
Overnight RMS (°C) & 0.35 & 0.30 & 0.33 \\
\bottomrule
\end{tabular}
\end{table}

The algorithm is robust to $\pm 25\%$ uncertainty in $\tau$.

\subsection{Sensitivity to Prediction Errors}

Rate-matching requires predicting $\dot{T}$. What if predictions are degraded?

\begin{table}[H]
\centering
\caption{Performance vs Prediction Quality}
\label{tab:pred_quality}
\begin{tabular}{l|cc}
\toprule
\textbf{Scenario} & \textbf{Sunset T RMS} & \textbf{Overnight T RMS} \\
\midrule
Perfect predictions & 0.00°C & 0.18°C \\
RF predictions (actual) & 0.69°C & 0.30°C \\
50\% degraded predictions & 0.95°C & 0.38°C \\
Temperature-only control & 0.80°C & 0.53°C \\
\bottomrule
\end{tabular}
\end{table}

Even with significantly degraded predictions, rate-matching outperforms temperature-only control.

\subsection{Worst-Case Nights}

The 95th percentile errors are:
\begin{itemize}
    \item Temperature-Only: 1.05°C overnight RMS
    \item Rate-Matching: 0.58°C overnight RMS (45\% better)
\end{itemize}

%=============================================================================
\section{Implementation Checklist}
%=============================================================================

\subsection{Required Components}

\begin{enumerate}
    \item \textbf{Prediction Models}: Random Forest models for $T_\text{sunset}$ and $\dot{T}_\text{sunset}$
    \begin{itemize}
        \item Train on historical data (odd calendar days)
        \item Retrain monthly with new data
        \item Input: temperature history, seasonality features
    \end{itemize}

    \item \textbf{Real-time Measurements}:
    \begin{itemize}
        \item Ambient temperature sensor (15-min resolution)
        \item Mirror temperature sensors (multiple locations)
    \end{itemize}

    \item \textbf{Control System Interface}:
    \begin{itemize}
        \item Setpoint update capability (every 15 minutes)
        \item Setpoint range: at least $\pm 5°C$ from ambient
    \end{itemize}
\end{enumerate}

\subsection{Operational Parameters}

\begin{table}[H]
\centering
\caption{Recommended Control Parameters}
\label{tab:params}
\begin{tabular}{ll}
\toprule
\textbf{Parameter} & \textbf{Value} \\
\midrule
Mirror thermal time constant ($\tau$) & 2.0 hours \\
Dome air time constant ($\tau_\text{dome}$) & 0.5 hours (30 min) \\
\textbf{Pre-sunset cold bias} & \textbf{0.3°C} \\
\textbf{Overnight cold bias} & \textbf{0.3°C} \\
Max setpoint rate ($R_\text{max}$) & 2.0°C/hour \\
Control update interval & 15 minutes \\
Rate estimation window & 1 hour \\
Pre-sunset phase start & 4 hours before sunset \\
\bottomrule
\end{tabular}
\end{table}

%=============================================================================
\section{Conclusion and Recommendations}
%=============================================================================

\subsection{Key Findings}

\begin{enumerate}
    \item \textbf{Rate-matching improves all metrics}: 13\% better temperature, 61\% better rate at sunset; 44\% better overnight tracking

    \item \textbf{Cold bias achieved}: 60\% of nights are slightly cold (vs 35\% with temperature-only control)

    \item \textbf{Robust to uncertainties}: Algorithm performs well even with imperfect predictions or time constant estimates

    \item \textbf{Simple implementation}: Only requires RF predictions for two quantities and standard setpoint control
\end{enumerate}

\subsection{Primary Recommendation}

\begin{center}
\fbox{\parbox{0.9\textwidth}{
\textbf{Dual-Setpoint Control Algorithm:}

\textbf{Pre-sunset} (dome closed):
\begin{equation*}
    T_\text{HVAC} = T_\text{predicted} - 0.3°C \quad \text{(FLAT, constant)}
\end{equation*}
Mirror naturally equilibrates to dome interior temperature.

\textbf{Overnight} (dome open):
\begin{equation*}
    T_\text{mirror,sp} = T_\text{ambient} + 2\text{h} \times \dot{T}_\text{ambient} - 0.3°C
\end{equation*}

Setpoint rate limited to $\pm 2°C$/hour.
}}
\end{center}

\subsection{Expected Performance}

Based on simulations across 769 nights:
\begin{itemize}
    \item \textbf{Overnight RMS Error: 0.44°C}
    \item \textbf{Mean Offset: $-0.26$°C} (a few tenths cooler than ambient, as desired)
    \item \textbf{Only 15\% of time warm} (significantly reduced)
    \item \textbf{68\% of time in ideal range} (mirror 0 to 0.5°C cooler than ambient)
\end{itemize}

\subsection{Future Enhancements}

Potential improvements to investigate:
\begin{itemize}
    \item Adaptive $\tau$ estimation from mirror response data
    \item Weather-conditional cold bias (more aggressive in humid conditions)
    \item Prediction confidence-weighted control (more aggressive when predictions are certain)
    \item Multi-zone control for temperature gradient management
\end{itemize}

%=============================================================================
\appendix
\section{Mathematical Derivation}
%=============================================================================

\subsection{Rate-Matching Setpoint Formula}

Given mirror dynamics (Eq.~\ref{eq:dynamics}):
\begin{equation}
    \frac{dT_m}{dt} = \frac{T_{sp} - T_m}{\tau}
\end{equation}

At steady state, $dT_m/dt = 0$ implies $T_m = T_{sp}$.

For rate-matching, we want $dT_m/dt = \dot{T}_\text{target}$ when $T_m = T_\text{target}$:
\begin{equation}
    \dot{T}_\text{target} = \frac{T_{sp} - T_\text{target}}{\tau}
\end{equation}

Solving for setpoint:
\begin{equation}
    T_{sp} = T_\text{target} + \tau \cdot \dot{T}_\text{target}
\end{equation}

This is the fundamental rate-matching formula. The cold bias term is added as an operational preference.

\subsection{Lead Compensation Interpretation}

The term $\tau \cdot \dot{T}_\text{target}$ acts as lead compensation:
\begin{itemize}
    \item If ambient is cooling at $-1°C$/h, we set the mirror 2°C colder
    \item The mirror then warms toward the setpoint at $+1°C$/h
    \item This warming rate matches the ambient cooling rate
    \item The mirror tracks ambient despite the thermal lag
\end{itemize}

\end{document}
